\scanpage{63}
\mspage{\uline{19}}
of the diagrams 17), gives a $\underline{\mathrm{Lim}}$ which is
represented likewise as a sum space, has little
chance of being homeomorphic to $X$ (for example if $X$ connected\ldots) !

The $V_{i,k}^j$ are germs of subspaces
of $X$, their only canonical \deleted{\unsure{relations of i-}} \inserted{morphisms}
between them are the inclusion morphisms
— under what condition then does one have
\[
(18)\qquad V_{i,k}^j\subset V_{i',k'}^{j'}
\]
for
\[
(19)\qquad i\leq j\leq k,\qquad i'\leq j'\unsure{\leq|\lneq}k'
\]
given ? What about the spaces
elementary
\[
X_i^*,\quad V_{i,j}^*,\quad V_{i,j}\qquad(j\text{ successor of }i)\ ?
\]
I shall assume for simplicity the $X_i^*\ne\varnothing$.
\inserted{(\deleted{\unsure{This|then}} $X_i\subset X_j\Longrightarrow i\leq j$ — for otherwise one would have $X_i^*=\unsure{\varnothing}$ precisely\ldots)}
There are thus no inclusion relations between
the $X_i^*$ for distinct indices, \deleted{\unsure{if for}} \inserted{If}
$V_{i,j}\subset V_{i',j'}$ \unsure{let us think|let us take|passing} to the ``points'' \inserted{of $X$''} of these germs
of spaces, one must have $X_i^*\subset X_{i'}^*$ whence $i=i'$
so if one assumes that the germ of an
$X_j$ ($j$ successor of $i$) along $X_i^*$ determines
$j$ already (which is quite \unsure{plausible}), one
will also have to have $j=j'$. Likewise, one sees that
if $V_{ij}^*\subset V_{i',j'}^*$, \unsure{their|their} sets of adherent points,
\deleted{\unsure{situated on|namely} $\overline{X_i^*}$ and $\overline{X_{i'}^*}$}, lying between $\overline{X_i^*}$ and $X_i$
resp $\overline{X_{i'}^*}$ and $X_{i'}$, are in an inclusion relation,
\margin{62}
